\section{弱型控制、全体高矩首项与限制型范数}\label{mom:section}

本节的作用是把光晕的顶层行为转化为归一化示性函数的高矩行为。这里的弱型估计只针对 Lebesgue 测度；它不是关于任意底测度的一致结论。记
\[
 q_\beta:=\frac{n}{n-\beta}.
\]
对一个几何基 $\mathcal B$ 和 $p>q_\beta$，定义限制型泛函
\begin{equation}
 \mathfrak F_{\beta,\mathcal B}(p):=
 \sup_{0<|E|<\infty}\frac1{|E|}
 \int_{\mathbb R^n}\left(\frac{M_{\beta,\mathcal B}\mathbf 1_E(x)}{|E|^{\beta/n}}\right)^p\,dx .
 \label{mom:functional}
\end{equation}
若 $r_p^{-1}=p^{-1}+\beta/n$，则 $1<r_p<n/\beta$，且 $\mathfrak F_{\beta,\mathcal B}(p)^{1/p}$ 正是此 Sobolev 线上限制于示性函数的 $L^{r_p}\to L^p$ 最优常数。

\begin{lemma}[有限中心方体的选择]\label{mom:selection}
设 $Q_j=x_j+[-r_j,r_j]^n$（$1\le j\le N$）为有限个中心方体。存在子族 $\{Q_{j_\nu}\}_{\nu=1}^M$，满足：
\begin{enumerate}
\item 对每个 $j$，存在 $\nu$ 使 $r_{j_\nu}\ge r_j$ 且 $x_j\in Q_{j_\nu}$；
\item $\sum_{\nu=1}^M\mathbf 1_{Q_{j_\nu}}\le2^n$ 几乎处处。
\end{enumerate}
\end{lemma}
\begin{proof}
按半边长递减排列。依次扫描方体；其中心不属于此前已选方体时才保留它。被舍弃的方体即满足第一项。

取一个不在所有坐标超平面 $\{y_k=(x_j)_k\}$ 上的点 $y$。对于包含 $y$ 的已选方体，按 $x_j-y$ 的坐标符号分为 $2^n$ 个正交象限。若同一象限内 $Q_i,Q_j$ 均包含 $y$ 且 $r_i\ge r_j$，则逐坐标有
\[
 |(x_i)_k-(x_j)_k|\le r_i,
\]
所以 $x_j\in Q_i$。在扫描中 $Q_j$ 本应被舍弃，矛盾。因此每个象限至多对应一个已选方体。除去零测的坐标超平面并得第二项。
\end{proof}

\begin{lemma}[紧集上的有限稳定见证]\label{mom:stable-witness}
设 $f\in L^1(\mathbb R^n)$，$\lambda>0$，且
\[
 \Omega_\lambda=\{x:M_{\beta,\mathcal Q_c}f(x)>\lambda\}.
\]
对每个紧集 $K\subset\Omega_\lambda$ 和 $0<\varepsilon<1$，存在有限个 $x_j\in K$、$r_j>0$ 及开集 $U_j\ni x_j$，使 $K\subset\bigcup_jU_j$，$U_j\subset x_j+(-\varepsilon r_j,\varepsilon r_j)^n$，且对每个 $z\in U_j$，方体 $z+[-r_j,r_j]^n$ 都严格见证 $M_{\beta,\mathcal Q_c}f(z)>\lambda$。
\end{lemma}
\begin{proof}
对每个 $x\in K$ 选择一个严格见证半边长 $r_x$。固定半边长的卷积平均关于中心连续，故该严格不等式在 $x$ 的某个开邻域内仍成立。再与 $x+(-\varepsilon r_x,\varepsilon r_x)^n$ 相交，最后对 $K$ 取有限子覆盖。
\end{proof}

\begin{theorem}[中心方体的分数弱型估计]\label{mom:centered-weak}
对 $f\in L^1(\mathbb R^n)$ 和 $\lambda>0$，
\begin{equation}
 \lambda\,|\{M_{\beta,\mathcal Q_c}f>\lambda\}|^{1/q_\beta}
 \le 2^{n/q_\beta}\|f\|_1=2^{n-\beta}\|f\|_1.
 \label{mom:weak}
\end{equation}
因此
\begin{equation}
 C_{\beta,\mathcal Q_c}(\alpha)\le 2^n\alpha^{-q_\beta}\qquad(0<\alpha<1).
 \label{mom:centered-tauberian-bound}
\end{equation}
\end{theorem}
\begin{proof}
令 $K$ 为严格水平集的紧子集，并应用引理~\ref{mom:stable-witness}。对所得有限见证族使用引理~\ref{mom:selection}。若 $z\in K$，则某个稳定见证的中心落在某个已选方体中，且其半边长不超过已选方体的半边长；所以
\[
 K\subset\bigcup_{\nu=1}^M\bigl(x_{j_\nu}+[-(1+\varepsilon)r_{j_\nu},(1+\varepsilon)r_{j_\nu}]^n\bigr).
\]
记 $a_\nu=\int_{Q_{j_\nu}}|f|$。严格见证给出
\[
 \lambda |Q_{j_\nu}|^{1/q_\beta}<a_\nu.
\]
于是由 $q_\beta>1$、Minkowski 不等式和引理~\ref{mom:selection}，
\[
 \begin{split}
 \lambda |K|^{1/q_\beta}
 &\le(1+\varepsilon)^{n/q_\beta}
 \Bigl(\sum_\nu a_\nu^{q_\beta}\Bigr)^{1/q_\beta}\\
 &\le(1+\varepsilon)^{n/q_\beta}
 \int |f(y)|\Bigl(\sum_\nu\mathbf1_{Q_{j_\nu}}(y)\Bigr)^{1/q_\beta}\,dy\\
 &\le(1+\varepsilon)^{n/q_\beta}2^{n/q_\beta}\|f\|_1.
 \end{split}
\]
令 $\varepsilon\downarrow0$，再用开集的 Lebesgue 内正则性，得到 \eqref{mom:weak}。令 $f=\mathbf1_E$、$\lambda=\alpha|E|^{\beta/n}$，并利用 $q_\beta\beta/n=q_\beta-1$，即得 \eqref{mom:centered-tauberian-bound}。
\end{proof}

\begin{remark}\label{mom:weak-scope}
常数 $2^{n-\beta}$ 来自有限中心方体的 $2^n$ 重叠界。定理没有断言该常数是 Lebesgue Tauberian 常数或任意测度弱型常数的最佳值。
\end{remark}

\begin{proposition}[全域低阈值支配]\label{mom:low-level}
存在仅依赖 $n$ 的常数 $K_n$，使
\[
 C_{\beta,\mathcal Q}(\alpha)\le3^n\alpha^{-q_\beta},\qquad
 C_{\beta,\mathcal Q_c}(\alpha)\le2^n\alpha^{-q_\beta},
\]
并且
\[
 C_{\beta,\mathcal R}(\alpha)
 \le K_n^2\alpha^{-2/(1-\beta/n)}.
\]
\end{proposition}
\begin{proof}
中心情形已经由定理~\ref{mom:centered-weak} 给出。以下证明非中心方体估计，以便记录低阈值支配实际使用的选择步骤。令
\[
 \Omega=H_{\beta,\mathcal Q}(\alpha,E),\qquad
 \lambda=\alpha|E|^{\beta/n}.
\]
严格见证在小的外扩下仍严格，故对每个紧集 $K\subset\Omega$ 可取有限个严格见证方体 $Q_1,\ldots,Q_N$，使 $K\subset\bigcup_j\operatorname{int}Q_j$。按边长递减贪心地保留与此前已选方体不相交的方体，记所得子族为 $Q_{j_1},\ldots,Q_{j_M}$。每个被舍弃的 $Q_j$ 与某个不小于它的已选方体相交，故逐坐标比较区间得到 $Q_j\subset3Q_{j_\nu}$。于是
\[
 |K|\le3^n\sum_{\nu=1}^M|Q_{j_\nu}|.
\]
令 $s_\nu=|E\cap Q_{j_\nu}|$。严格见证和 $q_\beta=(1-\beta/n)^{-1}$ 给出
\[
 |Q_{j_\nu}|<\lambda^{-q_\beta}s_\nu^{q_\beta}.
\]
所选方体两两不交，故 $\sum_\nu s_\nu\le|E|$；又 $q_\beta>1$，从而
\[
 \lambda^{-q_\beta}\sum_\nu s_\nu^{q_\beta}
 \le\lambda^{-q_\beta}|E|^{q_\beta}
 =\alpha^{-q_\beta}|E|.
\]
对紧 $K$ 取上确界即得第一个不等式。

对矩形，若 $s=|E\cap R|$，则
\[
 |R|^{\beta/n-1}s/|E|^{\beta/n}
 \le(s/|R|)^{1-\beta/n}.
\]
故相应光晕包含于经典强极大函数的水平集 $\{M_{0,\mathcal R}\mathbf1_E>\alpha^{1/(1-\beta/n)}\}$。经典强极大函数在 $L^2$ 上有仅依赖 $n$ 的范数界；这可由一维 Hardy--Littlewood 极大函数逐坐标迭代得到，亦见 de Guzm\'an \cite{deGuzman1975}。再用 Chebyshev 不等式即得结论。
\end{proof}

\begin{theorem}[固定方体的高矩首项]\label{mom:fixed-cube-moment}
设 $Q_0=[-1/2,1/2]^n$。由定理~\ref{tan:unit-cube-limit}，
\[
 |H_{\beta,\mathcal Q_c}(1-t,Q_0)|=|S_{n,c}|t^n+o(t^n)\quad(t\downarrow0),
\]
其中 $c=\min\{\beta,n-\beta\}$，$k=\lfloor2c\rfloor$，并且
\begin{equation}
 \lim_{p\to\infty}p^n\int_{\mathbb R^n}
 \bigl(M_{\beta,\mathcal Q_c}\mathbf1_{Q_0}\bigr)^p\,dx
 =n!|S_{n,c}|=
 \frac{n!\,2^n}{k!(2c)^{n-k}}.
 \label{mom:fixed-cube-asymptotic}
\end{equation}
\end{theorem}
\begin{proof}
层析公式和 $|Q_0|=1$ 给出
\[
 \int(M_{\beta,\mathcal Q_c}\mathbf1_{Q_0})^p
 =p\int_0^1\alpha^{p-1}|H_{\beta,\mathcal Q_c}(\alpha,Q_0)|\,d\alpha.
\]
在充分靠近 $1$ 的区间，中心端点上界给出 $|H_{\beta,\mathcal Q_c}(\alpha,Q_0)|\le A(1-\alpha)^n$；远离 $1$ 的部分由命题~\ref{mom:low-level} 指数衰减。故可在近端使用支配收敛，再用
\[
 p\int_0^1\alpha^{p-1}(1-\alpha)^n\,d\alpha=pB(p,n+1)\sim n!p^{-n}.
\]
这证明结论。
\end{proof}

\begin{theorem}[限制型高矩渐近]\label{mom:restricted-growth}
记 $L_{n,\beta}$ 如定理~\ref{gl:main}。当 $p\to\infty$ 时，
\[
 \mathfrak F_{\beta,\mathcal Q}(p)=1+O_{n,\beta}(p^{-1}),\qquad
 \mathfrak F_{\beta,\mathcal R}(p)=1+O_{n,\beta}(p^{-1}),
\]
且
\[
 \mathfrak F_{\beta,\mathcal Q_c}(p)
 =n!L_{n,\beta}p^{-n}+O_{n,\beta}(p^{-n-1}).
\]
中心公式的误差还在任意 $\beta\in J\Subset(0,n)$ 上一致（$p$ 充分大取决于 $n,J$）。
此外，非中心方体和强矩形满足
\[
 \liminf_{p\to\infty}p\bigl(\mathfrak F_{\beta,\mathcal Q}(p)-1\bigr)
 \ge\frac{2n}{n-\beta},
 \qquad
 \liminf_{p\to\infty}p\bigl(\mathfrak F_{\beta,\mathcal R}(p)-1\bigr)
 \ge\frac{2n^2}{n-\beta}.
\]
在 $n=1$，由一维精确公式，两个相应极限分别为
\[
 \lim_{p\to\infty}p\bigl(\mathfrak F_{\beta,\mathcal Q}(p)-1\bigr)=\frac{2}{1-\beta},
 \qquad
 \lim_{p\to\infty}p\mathfrak F_{\beta,\mathcal Q_c}(p)
 =\max\!\left\{\frac1\beta,\frac1{1-\beta}\right\}.
\]
\end{theorem}
\begin{proof}
对任意 $E$，自然顶层和层析公式给出
\[
 \frac1{|E|}\int\left(\frac{M_{\beta,\mathcal B}\mathbf1_E}{|E|^{\beta/n}}\right)^p
 =p\int_0^1\alpha^{p-1}
 \frac{|H_{\beta,\mathcal B}(\alpha,E)|}{|E|}\,d\alpha.
\]
对上确界使用端点光晕上界，并由命题~\ref{mom:low-level} 支配远离 $1$ 的区间。Beta 积分
\[
 p\int_0^1\alpha^{p-1}(1-\alpha)^j\,d\alpha\sim j!p^{-j}
\]
给出非中心的两个上界。中心情形使用定理~\ref{gl:main}：对每个输入，
近端层析积分至多
$L_{n,\beta}pB(p,n+1)+O_{n,J}(pB(p,n+2))$，
远端由弱型界一致指数衰减。下界取固定 $Q_0$ 并用其 $N=1$ 一致截断，得到同样的表达式。
由于 $pB(p,n+1)=n!\prod_{j=1}^n(p+j)^{-1}$、$pB(p,n+2)=O_n(p^{-n-1})$，
中心结论及紧参数一致性成立。这里仅将上确界放在积分外作上界，没有交换二者为等式。
为得到两个非中心下界，固定 $E=[0,1]^n$。实际测试方体包含给出
\[
 C_{\beta,\mathcal Q}(1-t)\ge
 \bigl(2(1-t)^{-1/(n-\beta)}-1\bigr)^n
 =1+\frac{2n}{n-\beta}t+O(t^2),
\]
以及
\[
 C_{\beta,\mathcal R}(1-t)\ge
 1+2n\bigl((1-t)^{-n/(n-\beta)}-1\bigr)
 =1+\frac{2n^2}{n-\beta}t+O(t^2).
\]
把这两个固定测试集下界代入层析公式，并使用
\[
 p\int_0^1\alpha^{p-1}(1-\alpha)\,d\alpha=\frac1{p+1},
\]
即得所列两个 \(\liminf\)。一维结论由一维精确常数在固定近端区间内达到，再作同一 Beta 积分得到。
\end{proof}

\begin{remark}\label{mom:global-gap}
定理~\ref{gl:main}闭合中心全体首项后，限制型泛函也具有匹配的全体首项。
下一系数仍不能仅由固定输入 $Q_0$ 的高矩展开推出，因为取上确界的输入可以随 $p$ 改变。
\end{remark}

\begin{corollary}[二维半阶的全体高矩次阶夹逼]\label{mom:pair-bracket}
当 $n=2$、$\beta=1$ 时，
\[
 0\le\liminf_{p\to\infty}p^3[\mathfrak F_{1,\mathcal Q_c}(p)-4p^{-2}]
 \le\limsup_{p\to\infty}p^3[\mathfrak F_{1,\mathcal Q_c}(p)-4p^{-2}]\le8.
\]
\end{corollary}
\begin{proof}
定理~\ref{pb:volume}对每个输入给出统一近端层析上界。由
$\mathcal U(1-t)=2t^2+(10/3)t^3+O(t^4)$ 及低水平弱型支配，
\[
 \mathfrak F_{1,\mathcal Q_c}(p)
 \le2pB(p,3)+\tfrac{10}{3}pB(p,4)+O(p^{-4})
 =4p^{-2}+8p^{-3}+O(p^{-4}).
\]
下界使用固定 $Q_0$ 的面积展开 $2t^2+2t^3+O(t^4)$，
而非把矩形上确界放进积分；它给出
$\mathfrak F_{1,\mathcal Q_c}(p)\ge2pB(p,3)+2pB(p,4)+O(p^{-4})
=4p^{-2}+O(p^{-4})$。两个不等式合并即得结论。
\end{proof}


\begin{corollary}[高矩层析包络的显式改进]\label{mom:sixth-envelope}
固定 $n=2$、$\beta=1$，令 $a_0=1/\sqrt2$。对 $p>2$ 定义
\[
 \mathcal B(p)=\frac{4p}{p-2}a_0^{p-2}
 +p\int_{a_0}^1\alpha^{p-1}\mathcal U(\alpha)\,d\alpha.
\]
则
\[
 \mathfrak F_{1,\mathcal Q_c}(p)\le\mathcal B(p)-\mathcal D(p),\qquad
 \mathcal D(p)=\frac p{768}\int_{99/100}^1\alpha^{p-1}(1-\alpha)^6\,d\alpha>0,
\]
且
\[
 \mathcal D(p)=\frac{15}{16}\prod_{j=1}^6(p+j)^{-1}
 +O\bigl((99/100)^p\bigr),\qquad
 \lim_{p\to\infty}p^6\mathcal D(p)=\frac{15}{16}.
\]
这里改进的是统一层析上包络，没有确定真实泛函的六阶系数。
\end{corollary}
\begin{proof}
由中心弱型估计~\eqref{mom:weak}，每个质量一输入的光晕面积至多
$4\alpha^{-2}$；在 $[a_0,1]$ 上改用定理~\ref{pb:volume}，并在
$[99/100,1]$ 上再减去定理~\ref{pb:sixth-gap}给出的 $(1-\alpha)^6/768$。
逐输入层析后对输入取上确界，即得所列不等式。
完整 Beta 积分为
\[
 \frac p{768}\int_0^1\alpha^{p-1}(1-\alpha)^6\,d\alpha
 =\frac{15}{16}\prod_{j=1}^6(p+j)^{-1},
\]
舍去的 $[0,99/100]$ 积分介于 $0$ 与 $(99/100)^p/768$ 之间，
故余项和极限成立。
\end{proof}


\begin{corollary}[抛物邻域对高矩包络的改进]\label{mom:parabolic-envelope}
仍记 $\mathcal B(p)$ 如推论~\ref{mom:sixth-envelope}。则对 $p>2$，
\[
 \mathfrak F_{1,\mathcal Q_c}(p)\le\mathcal B(p)-\mathcal D_*(p),\qquad
 \mathcal D_*(p)=2^{-24}p\int_{99/100}^1\alpha^{p-1}(1-\alpha)^{9/2}\,d\alpha,
\]
且
\[
 \lim_{p\to\infty}p^{9/2}\mathcal D_*(p)
 =2^{-24}\Gamma(11/2)=\frac{945\sqrt\pi}{2^{29}}.
\]
\end{corollary}
\begin{proof}
逐输入使用定理~\ref{pb:parabolic-gap}并重复前一推论的层析论证。
完整积分等于 $2^{-24}pB(p,11/2)$，截去的低端至多为 $2^{-24}(99/100)^p$。
由 Beta 积分的 Gamma 比值渐近即得极限。
这仍是统一上包络的减项，而非真实泛函的高阶系数。
\end{proof}


\begin{corollary}[高矩包络的四阶减项]\label{mom:quartic-envelope}
存在绝对 $\kappa,t_0>0$，使对 $p>2$，
\[
 \mathfrak F_{1,\mathcal Q_c}(p)\le\mathcal B(p)-\mathcal D_4(p),\qquad
 \mathcal D_4(p)=\kappa p\int_{1-t_0}^1\alpha^{p-1}(1-\alpha)^4\,d\alpha,
\]
且 $\lim_{p\to\infty}p^4\mathcal D_4(p)=24\kappa>0$。
这里可将 $t_0$ 缩小到 $t_0<1-1/\sqrt2$。
\end{corollary}
\begin{proof}
逐输入使用定理~\ref{pb:quartic-gap}，然后层析并对输入取上确界。
完整 Beta 积分为 $24\kappa\prod_{j=1}^4(p+j)^{-1}$，
截去部分至多为 $\kappa(1-t_0)^p$。
仍未据此确定真实泛函的四阶系数。
\end{proof}

\begin{corollary}[三阶亏损对高矩次项的严格改善]\label{mom:cubic-envelope}
取定理~\ref{ic:cubic-gap}中的 $\kappa,t_0$，并缩小 $t_0$ 使 $1-t_0>1/\sqrt2$。
则对 $p>2$，
\[
 \mathfrak F_{1,\mathcal Q_c}(p)\le\mathcal B(p)-\mathcal D_3(p),\qquad
 \mathcal D_3(p)=\kappa p\int_{1-t_0}^1\alpha^{p-1}(1-\alpha)^3\,d\alpha,
\]
且 $p^3\mathcal D_3(p)\to6\kappa>0$。特别地
\[
 0\le\liminf_{p\to\infty}p^3[\mathfrak F_{1,\mathcal Q_c}(p)-4p^{-2}]
 \le\limsup_{p\to\infty}p^3[\mathfrak F_{1,\mathcal Q_c}(p)-4p^{-2}]
 \le8-6\kappa<8.
\]
\end{corollary}
\begin{proof}
逐输入应用三阶亏损，再按层析公式积分并取上确界。
完整 Beta 积分为 $6\kappa\prod_{j=1}^3(p+j)^{-1}$，
截去的低端至多为 $\kappa(1-t_0)^p$。
结合推论~\ref{mom:pair-bracket}证明中的展开及固定 $Q_0$ 下界，即得次项夹逼。
这改善了真实泛函的统一上界，尚未给出其精确三阶系数。
\end{proof}
